% !TeX program = xelatex
\documentclass[
  a4paper,fontsize=11pt,parskip=half,headings=small,
  numbers=noenddot,bibliography=totoc
]{scrartcl}
\usepackage[margin=16mm,includeheadfoot,headheight=14pt,headsep=6mm,footskip=9mm]{geometry}
\usepackage{fontspec}
\usepackage{mathtools}
\usepackage{unicode-math}
\setmainfont{TeX Gyre Pagella}
\setsansfont{TeX Gyre Heros}[Scale=MatchLowercase]
\setmonofont{Latin Modern Mono}[Scale=MatchLowercase]
\setmathfont{TeX Gyre Pagella Math}
\usepackage[protrusion=true,expansion=false]{microtype}
\usepackage{tabularray}
\UseTblrLibrary{booktabs}
\usepackage{booktabs}
\usepackage{enumitem}
\usepackage{xcolor}
\usepackage{scrlayer-scrpage}
\usepackage[authoryear,round]{natbib}
\usepackage{xurl}
\usepackage[
  unicode,colorlinks=true,linkcolor=blue!50!black,
  citecolor=blue!50!black,urlcolor=blue!50!black,
  pdftitle={Multiple-Test Admissions: A Joint Analysis of Two Grady Towers Letters},
  pdfsubject={Bivariate normal overlap, eligibility rules, and corrected conditional probabilities},
  pdfauthor={Prepared with OpenAI Codex},
  pdfkeywords={IQ, admissions, Grady Towers, bivariate normal, phi correlation, selection}
]{hyperref}
\usepackage{bookmark}
\usepackage[nameinlink,noabbrev]{cleveref}
\DeclareMathOperator{\Var}{Var}
\DeclareMathOperator{\Corr}{Corr}
\newcommand{\sfN}{\overline{\Phi}}
\setlength{\emergencystretch}{1.5em}
\setlist[itemize]{leftmargin=*,itemsep=2pt,topsep=3pt}
\setkomafont{disposition}{\normalfont\sffamily\bfseries}
\setkomafont{pageheadfoot}{\normalfont\sffamily\small}
\clearpairofpagestyles
\ihead{Multiple-test admissions}
\ohead{Joint analysis and corrected article}
\cfoot*{\pagemark}
\pagestyle{scrheadings}
\setlength{\bibsep}{5pt}
\title{Multiple-Test Admissions}
\subtitle{A joint analysis and corrected version of two Grady Towers letters}
\author{}
\date{1 November 2026}

\begin{document}
\maketitle
\thispagestyle{plain.scrheadings}
\vspace{-1.1em}
\begin{abstract}
Grady Towers's first letter gives incorrect overlap probabilities for
correlated IQ tests; his follow-up explicitly retracts that calculation and
introduces a substantially better numerical method. This report assesses
the two together. At correlation 0.7, a bivariate-normal model with two
top-2\% cutoffs gives 3.3724\% eligible through either test and 31.3814\%
repeat qualification among passers of one specified test. For top-2\% and
top-0.1\% cutoffs on different tests, 65.8589\% of higher-threshold passers
also pass the lower-threshold test. The follow-up's remaining numerical
differences arise chiefly from grid boundaries, rounded cutoffs, and
inconsistent marginal denominators. Neither letter establishes actual
membership proportions or a latent-ability interpretation. A corrected
combined article and reproducible calculations are provided.
\end{abstract}

\noindent\textbf{AI disclosure.} This report was substantially generated and revised
with OpenAI Codex and collaborating AI agents, including its analysis,
replacement text, code and verification workflow. Provenance details, human
involvement and limitations are disclosed in \cref{sec:ai-disclosure}.


\section{Joint assessment of the two sources}\label{sec:audit}

\subsection{Read the second letter as a retraction}

The sources are \emph{On using multiple tests for high IQ society admissions}
\citep{towersOriginal} and its \emph{Followup} \citep{towersFollowup}, both
hosted by Darryl Miyaguchi. The first is undated in the supplied page. The
second identifies its underlying letter as dated 12 January 2000 and says
that it was published posthumously with permission. The supplied Wayback
captures are dated 9 March 2026 and 5 May 2022, respectively; these are
archive dates, not the letters' publication dates.

The follow-up expressly withdraws the earlier calculation. It should not
be presented as independent support for the first letter's dramatic
conclusion. Towers's useful qualitative point survives: accepting a pass
on either of two imperfectly overlapping tests generally enlarges the
eligible group relative to either test alone. The numerical examples and
the interpretation of eligibility need correction.

This report examines the mathematical scenarios in those historical
letters. It does not establish the admission policies, test correlations,
or membership composition of any named society, either historically or
today. The incomplete 1960 reference in the first letter can be identified
as \citet{zelen1960}; the calculations here were performed independently.

\subsection{The first letter's principal numerical claims fail}

\Cref{tab:original-audit} compares the first letter's principal figures with
recalculations under its intended Gaussian scenario. In the mixed case,
\(A\) means passing the top-2\% cutoff on test \(X\), and \(B\) means
passing the top-0.1\% cutoff on a different test \(Y\). All corrected entries
assume population score correlation \(\rho=0.7\).

\begin{table}[htbp]
\centering
\caption{Selected claims in the original letter and corrected model results.
All displayed probabilities are percentages, with the denominator specified.}
\label{tab:original-audit}
\begin{tblr}{colspec={X[2.6,l] X[0.8,r] X[1,r]},row{1}={font=\bfseries},rowsep=4pt}
\toprule
Quantity & Original & Corrected \\
\midrule
Both top-2\% tests: population & about 0.2 & 0.627627 \\
Either top-2\% test: population & 3.83 & 3.372373 \\
Second top-2\% pass, given first pass & 10.27 & 31.381370 \\
Mixed case: either pass, population & 2.11 & 2.034141 \\
Mixed case: both pass, population & about 0.014 & 0.065859 \\
Mixed case: higher cutoff given lower, \(P(B\mid A)\) & 0.69 & 3.292944 \\
Mixed case: lower cutoff given higher, \(P(A\mid B)\) & 13.48 & 65.858880 \\
\bottomrule
\end{tblr}
\end{table}

One error is visible without a distributional model: the mixed-case union
cannot be 2.11\% when its marginal probabilities sum to only 2.10\%.
In general, \(P(A\cup B)\leq P(A)+P(B)\). The corrected mixed-case failure
rate on the lower-threshold test among higher-threshold passers is
34.1411\%, rather than the original's approximately 86\%.

For the equal top-2\% case, passing \(X\) alone has probability 1.3724\%,
and passing \(Y\) alone has the same probability. Passing \emph{exactly one}
therefore has probability 2.7447\%. These are different quantities; wording
such as ``one test but not the other'' should specify which is intended.

\subsection{What the follow-up fixes, and what remains}

The follow-up replaces an approximation involving the phi coefficient with
numerical summation of a bivariate-normal density. That is a substantial
correction. Reimplementing its supplied BASIC grid closely reproduces its
published joint probabilities. The remaining issues are:

\begin{itemize}
\item The equal-weight half-IQ-point grid samples the cutoff itself, giving
approximately the boundary of a midpoint cell a quarter point below that
cutoff. It does not integrate exactly from the stated continuous cutoff.
\item The integer IQ cutoffs 133, 150, and 164 do not have exactly the
nominal tail probabilities used as denominators in the follow-up.
\item A conditional proportion among passers of one specified test is not
automatically a proportion among society members, or even among all people
eligible through either of two tests.
\end{itemize}

\Cref{sec:grid} separates those effects. The absent scale and cell-area
factors in the density ratio cancel; they are not fatal normalization
errors. The image of the follow-up's double integral was also inspected,
in addition to its HTML text and BASIC listing.

The follow-up recounts a top-5\% example attributed to page 306 of Cattell's
\emph{The Scientific Analysis of Personality}. Under the model used here,
two such tests at \(\rho=0.7\) have joint probability 1.95993\%, rather than
1\%. This checks Towers's recollection as written. Cattell's cited page
was not independently inspected, so no error is attributed to Cattell
himself.

\section{Corrected combined article: admissions through correlated tests}
\label{sec:corrected}

An admission rule and a test percentile answer different questions. A test
cutoff may be exceeded by 2\% of a reference population. A rule accepting
either of two such cutoffs concerns the union of two sets of people, whose
size depends on their overlap. Requiring both cutoffs concerns their
intersection. Neither rule, by itself, identifies the people with the
highest values of a perfectly measured underlying ability.

\subsection{Specify a joint model and distinguish the correlations}

Let \(X,Y\) denote standardized scores in the same reference population,
with a bivariate-normal distribution:
\begin{equation}
\begin{pmatrix}X\\Y\end{pmatrix}
\sim N_2\!\left[
\begin{pmatrix}0\\0\end{pmatrix},
\begin{pmatrix}1&\rho\\\rho&1\end{pmatrix}\right],
\qquad -1<\rho<1.
\label{eq:model}
\end{equation}
Here \(\rho\) is the Pearson correlation of the continuous scores. It is
not the reliability of either test and is not generally the correlation of
the binary pass/fail indicators. Normal marginal distributions and a
correlation alone do not establish \cref{eq:model}; the joint distribution
is an additional assumption, especially consequential in the tails.

For desired marginal pass probabilities \(p_A,p_B\), use
\begin{equation}
a=\Phi^{-1}(1-p_A),\qquad b=\Phi^{-1}(1-p_B),
\qquad A=\{X>a\},\quad B=\{Y>b\}.
\label{eq:cutoffs}
\end{equation}
The joint pass probability follows by conditioning on \(X\):
\begin{equation}
\boxed{J=P(A\cap B)=\int_a^\infty
\varphi(x)\,\sfN\!\left(\frac{b-\rho x}{\sqrt{1-\rho^2}}\right)\,dx.}
\label{eq:joint}
\end{equation}
Here \(\varphi\) is the standard-normal density and
\(\sfN=1-\Phi\) its upper-tail function. The remaining quantities are
\begin{align}
P(A\cup B)&=p_A+p_B-J,\label{eq:union}\\
P(B\mid A)&=J/p_A, & P(A\mid B)&=J/p_B.\label{eq:conditionals}
\end{align}
The denominators determine the question being answered.

For binary indicators \(1_A,1_B\), the phi coefficient is
\begin{equation}
\phi_{\mathrm{binary}}=
\frac{J-p_Ap_B}{\sqrt{p_A(1-p_A)p_B(1-p_B)}}.
\label{eq:phi}
\end{equation}
This differs from the correlation of the underlying continuous Gaussian
variables \citep{ekstrom2011}. For equal top-2\% thresholds at
\(\rho=0.7\), \(\phi_{\mathrm{binary}}\approx0.29981\), rather than 0.7.
The follow-up attributes its earlier error to a phi/product-moment
approximation. Without the first algorithm, its exact computational
mistake cannot be reconstructed; substituting one correlation for the
other is nevertheless unjustified.

\subsection{Two tests with the same nominal cutoff}

For two top-2\% cutoffs and \(\rho=0.7\), \cref{eq:joint} gives
\begin{equation}
J=0.0062762740205,\qquad
P(A\cup B)=0.0337237259795.
\label{eq:equal-results}
\end{equation}
Thus 0.627627\% of the population passes both tests, and 3.372373\%
passes at least one. Among people who pass a specified first test,
\begin{equation}
P(B\mid A)=\frac{0.0062762740205}{0.02}
\approx31.3814\%.
\label{eq:repeat-pass}
\end{equation}
This statement concerns repeat qualification, not test reliability and
not an unconditional probability of joining a society.

If eligibility instead means passing \emph{either} of the two tests, the
fraction of that eligible union passing both is
\begin{equation}
P(A\cap B\mid A\cup B)=\frac{J}{2p-J}
\approx18.6109\%,\qquad p=0.02.
\label{eq:union-denominator}
\end{equation}
The 31.3814\% and 18.6109\% figures are both correct under the model;
they use different denominators. Actual members may be a further selected
subset of either group.

The same calculation at rarer equal cutoffs yields
\cref{tab:equal-cutoffs}. The cutoffs are exact nominal percentiles rather
than rounded IQ points. IQ-style values in the table use mean 100 and
standard deviation 16; they are scale conversions.

\begin{table}[htbp]
\centering
\caption{Equal nominal cutoffs at \(\rho=0.7\): population overlap and
repeat qualification among passers of one specified test.}
\label{tab:equal-cutoffs}
\begin{tblr}{colspec={X[1.1,l] X[0.8,r] X[1,r] X[1.1,r] X[1.1,r]},
row{1}={font=\bfseries},rowsep=4pt}
\toprule
Marginal tail on each test & \(z\) cutoff & SD-16 cutoff
 & Both: population \% & Second pass given first: \% \\
\midrule
2\% & 2.053749 & 132.85998 & 0.627627 & 31.381370 \\
0.1\% & 3.090232 & 149.44372 & 0.0159589 & 15.958919 \\
\(1/33{,}000\) & 4.010439 & 164.16702 & 0.000231250 & 7.631243 \\
\bottomrule
\end{tblr}
\end{table}

\subsection{A lower cutoff on one test and a higher cutoff on another}

Let \(p_A=0.02\) and \(p_B=0.001\), still with \(\rho=0.7\).
The joint probability is
\begin{equation}
J=0.0006585888039.
\label{eq:mixed-joint}
\end{equation}
Consequently,
\begin{align}
P(B\mid A)&=3.292944\%, &
P(A\mid B)&=65.858880\%,\label{eq:mixed-conditionals}\\
P(A^c\mid B)&=34.141120\%, &
P(A\cup B)&=2.034141\%.\label{eq:mixed-remainder}
\end{align}
The asymmetry is expected: the higher-threshold passers are a much
smaller group. These two events concern \emph{different tests}. If both
cutoffs were applied to the same score, the higher-cutoff event would be
contained in the lower-cutoff event and the corresponding conditional
probability would be 100\%.

The result is sensitive to the assumed correlation. For these same mixed
cutoffs, \(P(A\mid B)\) is approximately 33.59\% at \(\rho=0.5\),
65.86\% at 0.7, and 97.99\% at 0.9. A nominal 0.7 is an illustrative
parameter, not a universal empirical property of IQ tests. Population,
norms, age, constructs, administration, and selection all matter.

These conditional model results cannot determine whether someone who
passed a particular magazine test could join a particular society.
That requires evidence about the actual instruments and the society's
accepted routes. Nor does the model establish the earlier claim about the
prior memberships of most high-scoring society members.

\subsection{Eligibility size, rarity equivalents, and calibrated rules}

The population fraction admitted by a rule can be expressed as a normal
rarity equivalent. Under the equal top-2\% example, the OR and AND fractions
map to approximately 129.26 and 139.94, respectively, on an SD-16 scale:
\begin{equation}
I_{\mathrm{rarity}}=100+16\Phi^{-1}(1-p_{\mathrm{eligible}}).
\label{eq:rarity}
\end{equation}
This only matches a selected fraction. It does not show that the rules
select the same people as a single test at that cutoff, establish
``super-high reliability,'' or identify a particular level of latent
intelligence. The first letter's corresponding AND rarity equivalent
of about 146 follows from its erroneous overlap probability.

Multiple tests do not make a 2\% overall eligibility rule impossible.
Within the stated model, one can choose the common cutoff \(a_*\) to solve
\begin{equation}
2\sfN(a_*)-J(a_*,a_*;0.7)=0.02.
\label{eq:calibration}
\end{equation}
The solution is approximately \(a_*=2.270055\), or 136.32088 on an SD-16
scale. Each test then has marginal pass rate approximately 1.16021\%,
while the two-test OR rule has overall rate 2\%. This is an example of
calibrating a \emph{rule}; using more than one test is not inherently
incompatible with a specified eligible fraction.

\subsection{Interpret the population model and the testing process}

A useful report separates three questions: who exceeds a cutoff on one
instrument; who is eligible under a specified admission procedure; and
what the evidence supports about an individual's abilities. Agreement
on one does not settle the others. Score interpretations and reference
populations need explicit support
\citep[Standards~1.1--1.4 and 5.8--5.9]{standards2014}.

An eligibility probability refers to the designated population before
application or membership selection. Correlations measured among applicants
or members need not equal the population correlation used in
\cref{eq:model}. Selection can change the relationship, and adjustments
depend on the selection mechanism \citep{sackett2000}.

A procedure accepting the best result from repeated attempts also differs
from a fixed two-test procedure. For \(m\) pass events, independently of
any Gaussian model,
\begin{equation}
\max_i P(A_i)\leq P\!\left(\bigcup_{i=1}^{m}A_i\right)
\leq\min\!\left(1,\sum_{i=1}^{m}P(A_i)\right).
\label{eq:bounds}
\end{equation}
If attempts were independent with identical pass probability \(p\), the
union would be \(1-(1-p)^m\). Actual retests usually require a model of
dependence and administration history; practice and coaching can also
change score distributions \citep{hausknecht2007}. An analysis should
calibrate the procedure actually used, including its attempts, stopping
rule, and accepted instruments.

Towers was right to ask how multiple admission routes affect eligibility,
and right to withdraw a faulty calculation. A defensible replacement
states the rule, the reference population, the joint-distribution model,
the thresholds, and each conditional denominator. Its conclusions remain
conditional on those specifications.

\appendix
\section{Numerical audit of the follow-up}\label{sec:grid}

The BASIC listing uses \(\rho=0.7\), means 100, standard deviations 16,
and grid nodes from \(Q\) through 199 in steps of \(h=0.5\), with the
full-field denominator starting at 1. In mathematical notation its
probability estimate is
\begin{equation}
\widehat J_Q=
\frac{\displaystyle\sum_{x\in\{Q,Q+h,\ldots,199\}}
\sum_{y\in\{Q,Q+h,\ldots,199\}}g(x,y)}
{\displaystyle\sum_{x\in\{1,1+h,\ldots,199\}}
\sum_{y\in\{1,1+h,\ldots,199\}}g(x,y)},
\label{eq:grid}
\end{equation}
where \(g\) is proportional to the Gaussian density on the IQ-point
coordinates. The full grid has \(397^2=157{,}609\) nodes. A constant
factor \(1/(16\cdot16)\) and a cell area \(h^2\) would cancel between the
two sums. The corresponding absent scale factor also cancels in the
displayed double-integral ratio in the follow-up.

\begin{table}[htbp]
\centering
\caption{Reproduction of the follow-up's grid estimates. Entries are
probabilities, not percentages. The last column integrates from the actual
continuous cutoff \(Q\), on both tests.}
\label{tab:grid}
\begin{tblr}{colspec={X[0.45,l] X[1.4,r] X[1.4,r] X[1.4,r]},
row{1}={font=\bfseries},rowsep=4pt}
\toprule
\(Q\) & Published grid ratio & Reproduced grid ratio & Continuous joint probability \\
\midrule
133 & 0.0064061692 & 0.0064061681 & 0.0061136512 \\
150 & 0.00014748463 & 0.00014748460 & 0.00013830591 \\
164 & \(2.6393845\times10^{-6}\) & \(2.6393840\times10^{-6}\)
    & \(2.4390159\times10^{-6}\) \\
\bottomrule
\end{tblr}
\end{table}

The reproduced and published grid values agree to relative differences of
about \(2\times10^{-7}\), but not to every printed digit. The exact cause
of those small differences is not established. They are far smaller than
the difference between the grid estimate and the continuous probability
at the intended cutoff.

Equal weighting at the boundary includes a full contribution from the
node at \(Q\). Viewed as midpoint quadrature, its first cell begins at
\(Q-h/2=Q-0.25\). The effective lower boundary is therefore shifted
downward. The finite upper boundary and ordinary quadrature error add
smaller discrepancies in these examples. If IQ scores were explicitly
rounded to integers before classification, a separate rounding rule
would be needed; a quarter-point grid shift cannot simply be assumed to
be that rule.

For continuous scores, the actual marginal exceedance probabilities at
the stated integer cutoffs are those in \cref{tab:integer-cutoffs}.
Dividing the continuous joint tail by its matching marginal gives a
consistent conditional probability.

\begin{table}[htbp]
\centering
\caption{Consistent continuous calculations at the follow-up's integer
SD-16 cutoffs, with \(\rho=0.7\). All probability columns are percentages.}
\label{tab:integer-cutoffs}
\begin{tblr}{colspec={X[0.6,l] X[1.3,r] X[1.3,r] X[1.3,r]},
row{1}={font=\bfseries},rowsep=4pt}
\toprule
Cutoff & Marginal pass rate & Both: population & Second pass given first \\
\midrule
133 & 1.95800788 & 0.61136512 & 31.223834 \\
150 & 0.08890253 & 0.01383059 & 15.557028 \\
164 & 0.00316712 & 0.00024390 & 7.701043 \\
\bottomrule
\end{tblr}
\end{table}

The follow-up divides its grid values by 0.02, 0.001, and 0.000030303.
Those are approximations to desired nominal tails, not matching marginal
tails of the specific grid or continuous integer cutoffs. Its
approximately 32\%, 15\%, and 9\% summaries consequently mix conventions.
For exact nominal percentiles, use \cref{tab:equal-cutoffs}; for actual
continuous integer cutoffs, use \cref{tab:integer-cutoffs}. Using the
grid's own marginal denominator gives approximately 31.51\%, 15.74\%,
and 7.81\%, a third, internally defined approximation.

\section{Reproducibility and numerical checks}\label{sec:reproduction}

The accompanying script evaluates \cref{eq:joint} with adaptive
one-dimensional quadrature and directly uses inverse survival functions
for the thresholds. SciPy documents both its normal survival functions
and its integration routine \citep{scipyNorm,scipyQuad}.

An independent integration identity supplies a cross-check:
\begin{equation}
J(a,b;\rho)=\sfN(a)\sfN(b)+\int_0^\rho
\frac{\exp\!\left[-\dfrac{a^2-2tab+b^2}{2(1-t^2)}\right]}
{2\pi\sqrt{1-t^2}}\,dt.
\label{eq:rho-integral}
\end{equation}
This follows by integrating the Gaussian probability's derivative with
respect to its correlation, starting at independence. Checks also cover
symmetry, probability bounds, the original mixed-case union violation,
the binary phi coefficient, the calibrated OR threshold, and reproduction
of the follow-up's grid.

The CSV files contain unrounded model calculations to complement the
rounded display tables. Run \texttt{make check} from this analysis folder;
\texttt{make pdf} builds the report with XeLaTeX. Numerical consistency
does not establish the empirical truth of the Gaussian assumptions or
the chosen correlation.

\section{Expanded verification results and limits}\label{sec:verification}

The original proof selected only \texttt{Admission\_Core}. The expanded
implementation also puts \texttt{Normal\_Model}, \texttt{Towers\_Engine},
and the \texttt{Towers\_Demo} main procedure in SPARK, apart from small
explicit numerical completions. Parsing and CSV/file output are isolated
in \texttt{Towers\_IO}, with external state modeled in its SPARK interface.
The proof project now enumerates all thirteen source files, including the
trusted I/O completion. Turning an aspect on is a language-coverage change;
the fresh native proof results below provide the separate verification evidence.

The \href{https://github.com/AndersH3/composite_iq_checked_claimes/blob/main/verification/evidence/ghost-boundaries-computer10-20261101/REPORT.md}{Ghost-boundary computer10 campaign}
was completed on 1 November 2026. It records exact commands, tool versions, source
hashes, native reports and raw coverage, checksums, numerical comparisons,
and Python results. Earlier 31-check and 110-test results remain immutable
historical snapshots; they are not totals for this expanded implementation.

\subsection{SPARK scope and results}

The fresh five-prover portfolio discharged \textbf{498 checks:
123 flow checks and 375 prover checks}, with zero unproved
and zero justified checks. CVC5, Z3, Alt-Ergo, classic COLIBRI, and Yices were
requested together. A completed portfolio does not imply that every solver
completed every independent obligation. The deliberately incorrect union
control was rejected with an unproved functional contract. All invocations
use fresh source copies and object directories. The matching toolchain
reports GNAT 13.2.0 and GNATprove banner \texttt{0.0w}; its executable hashes
are recorded in the manifest.

Native coverage records 90 \texttt{all}, 16 \texttt{spec},
and zero \texttt{no} entities. The statuses describe analyzed bodies and
trusted specifications, not successful proofs by themselves. Whole off-mode
units can emit empty native entity arrays, so zero \texttt{no} entries do
not mean that no code is excluded. The explicit earlier-source baseline
has 11 \texttt{all} entities; it records flow/coverage only and cannot count
as a new proof. The auditor separately checks native proof totals, complete
metadata, and coverage of each required driver and numerical body.

The original core contracts still establish classification into the four
event cells, exactly one increment per observation, bounded marginal and
union counts, and exact integer/rational inclusion-exclusion, complements,
and conditional fractions. Expanded SPARK code covers inverse-tail
bisection, Simpson summation and grid loops, the exact Park--Miller state
transition and uniform conversion, Box--Muller composition, calculations,
simulation/count loops, Wilson arithmetic, result assembly, and the main
procedure. Joint-tail contracts require \(0\leq\rho\leq0.95\), positive even
panel counts, and \(\lvert A\rvert,\lvert B\rvert\leq10^{100}\); inversion
requires a probability strictly inside the endpoint survival bracket.
The main cases satisfy these implementation domains.
Small SPARK helpers and proved loop invariants bound accumulation
without importing proof axioms or changing the numerical summation order.

\subsection{Checked interfaces and numerical limitations}

Small \texttt{Normal\_Model} completions remain off for square root,
normal survival via foreign \texttt{erfc}, logarithm and exponential ranges,
exponent and quadrature node validation, probability bounds, and grid normalization.
The engine's checked completions validate quantization error and the
joint, union, and conditional relationships needed by the exact core. Each
promised normal-return bound or relationship is checked before returning;
an invalid value raises an error rather than being clamped or changed.
\texttt{Towers\_IO}'s parsing, formatting, filesystem operations, and
exceptions remain outside the proof.

Analysis uses these checked interfaces and the installed Ada
elementary-function specifications, whose contracts support SPARK but
provide limited numerical information. A successful caller proof does not
prove the implementations, that a numerical guard can never reject, or
that I/O succeeds. There are no project \texttt{pragma Assume} statements,
imported proof axioms, or proof suppressions. Zero explicit assumptions or
suppressions in native metadata does not remove trust in off-mode
completions and standard-library contracts. No proof of inversion accuracy,
survival-function monotonicity, quadrature error, or a Gaussian distribution
is claimed. The 10,000/20,000-panel comparison remains a numerical refinement
check, not a certified integration-error bound.

The floating-to-core conversion rounds probabilities to units of
\(10^{-15}\) and checks reconstruction error against
\(\delta=0.51\times10^{-15}\). For a conditional probability \(j/p\),
the resulting ratio change is bounded by
\(\delta(1+j/p)/(p-\delta)\), for \(p>\delta\). At the smallest nominal
marginal \(p=1/33{,}000\), this is at most approximately
\(3.37\times10^{-11}\); the checked changes are smaller. This accounts
only for quantization. Raw floating union and conditional values are
retained alongside the exact-core results.

\subsection{Ghost models and excluded implementation refinement}

The implemented Ghost package \texttt{Towers\_Boundary\_Models} specifies
the accepted normal-return
bounds and unchanged-value relationships of all eight numerical and four
engine completions. It also observes the actual typed sine and cosine
results. Their \([-1,1]\) interface model uses the installed library
specification; analytic accuracy, periodicity, and backend correctness
remain unproved. Each excluded completion calls a Ghost observer after
its ordinary checks. Original ordinary postconditions and concrete guards
remain in place when Ghost code is disabled.

\texttt{Towers\_IO\_Models} defines bounded decimal digit accumulation,
lexical field predicates, case labels, report validity, nominal layouts,
and an abstract receipt. Its decimal model defines a restricted unsigned
digit grammar using exact integer accumulation; its proved postconditions
establish bounds and nonempty recognition. Runtime observations compare
recognized inputs with Ada's parsed value. Other accepted Ada option forms
retain their original parser and a bounded-value observation; universal
parser equivalence is not proved. Field predicates do not certify numeric
round trips, formatting accuracy, or complete CSV syntax.

The Ghost receipt in the main procedure records returned options, the
requested data snapshot, four requested files and twelve requested data
rows, followed by an observed normal return. Its state-transition and
payload-preservation contracts are analyzed in SPARK. These quantities
are requests and observations, not proofs that files exist, contain those
bytes, were closed successfully, or persist on disk. The implementation's
refinement of the model and its failure or exception semantics remain
explicit obligations.

The audited register maps all native specification-only entities, the
foreign import, and the four private I/O helpers, including the nested
writer, to actual source-bound Ghost uses and analyzed model entities.
Separate cached sine/cosine assignments and their Ghost observations are
also checked. External library units and global unproved domains are
inventoried. Coverage of the register does not prove adequacy of its
models or correctness of the excluded implementations. A deliberately
wrong Ghost event counter is rejected alongside the original union control.

Fresh Ghost-enabled and Ghost-disabled builds keep ordinary assertions
enabled. Under both policies the million-pair calculations and simulations
reproduce all four original CSVs byte for byte, with validation and exact
replay passing in normal and optimized Python. Representative existing
Ada option forms, including signed, whitespace, underscore, based-literal
and exponent forms, also pass. These are finite regression checks;
SPARK's Ghost rules and the unproved excluded-code boundary remain distinct.
See \href{https://docs.adacore.com/live/wave/spark2014/html/spark2014_ug/en/source/specification_features.html}{AdaCore's specification guide}
for the modeling and erasure semantics.

\subsection{Numerical regression, Deal, and evidence audits}

A fresh Ada build reproduces all four retained CSVs byte for byte.
Independent SciPy validation and Python replay pass normally and under
\texttt{python -O}. One million shared correlated pairs from seed
\texttt{20261030} supply the four cases; all 16 event-cell counts match
exactly. Minimum valid trial and boundary-seed runs also succeed, and
invalid options are rejected. Replay checks reproducibility and agreement,
not random independence. Wilson intervals describe simulation uncertainty
under the model; rare events at this sample size cannot establish extreme-tail
psychometric frequencies.

All current Python files and named functions, including tests, fixtures,
and evidence tools, carry Deal contracts. The final repository-wide
inventory has 24 files, 213 named functions, and 607 Deal decorators with
zero gaps. All 23 positive files have zero lint diagnostics, and 209 tests
pass in each execution mode. Coverage, clean positive
static lint, the exact expected \texttt{DEL012} negative control,
runtime/property tests, numerical scripts, and historical-source audits
pass in normal and optimized Python. CLI entry points explicitly enable
Deal under optimization. Deal Solver proves the two pure proof-acceptance
gate postconditions and rejects the incorrect symbolic control. These
focused results do not prove the full Python programs, SciPy, or subprocess
handling; clean lint and finite tests are not universal proofs.

From a full clone, install the pinned dependencies and run
\texttt{make check}. Audit the current campaign with
\mbox{\texttt{python verification/verify\_evidence.py}}.
The auditor requires both production projects,
complete native coverage metadata, required SPARK-on bodies, consistent
native totals, source identities, checksums, and rejection of the negative
control. Historical Ada campaigns require an explicit matching old source
tree; changed Ada cannot silently fall back to history. Historical Python
identities may be resolved through verified Git blobs. Audits check record
integrity and identity; they do not rerun the provers or independently
validate SMT proof objects.

None of these software checks establishes that \(\rho=0.7\) applies to
actual tests, that empirical norms and tails are accurate, or that model
eligibility probabilities describe real membership selection or latent
ability. Those empirical and modeling questions remain as identified in
the main analysis.

\clearpage
\section{AI disclosure and provenance}\label{sec:ai-disclosure}

This report was substantially generated and revised using OpenAI Codex
and collaborating AI sub-agents. The session identifies Codex as based on
the GPT-6 family. An exact serving checkpoint, per-agent model identifiers,
and generation settings were not preserved as versioned provenance; no
more precise model attribution is claimed. This disclosure applies to the
analysis, corrected replacement text, project-authored software, and earlier revisions,
as well as the present additions.

AI contributions included interpreting the supplied articles, drafting and
revising the mathematical exposition and statistical discussion, selecting
and checking numerical examples, researching references, designing and
editing the Python/Deal and Ada/SPARK programs, ghost models, tests and
evidence auditors, orchestrating verification, typesetting and inspecting
the PDFs, and preparing repository commits. The human requester supplied
the source articles, specified the objectives, formats and verification
requirements, requested revisions, and authorized repository publication.
The recorded audits and parallel reviews were performed by AI agents and
automated tools. Independent human peer review or professional sign-off
is not recorded and is not claimed. Repository ownership and Git commit
attribution do not identify passages as human-authored.

The original authors remain credited as sources. The replacement text is
an AI-generated analytical reconstruction prepared for this project;
no endorsement by the original authors, their hosts, or cited organizations
is claimed. Sources were consulted at different levels, including article
text, software documentation and bibliographic metadata. Per-reference
notes identify relevant access limits; full-text review of every cited
work is not claimed.

Calculations, tests and prover invocations were actually executed on
\texttt{computer10}; the retained commands, native outputs, source hashes,
tool versions and checksums provide the verification evidence. These are
machine-checker results under the documented contracts and assumptions,
not evidence supplied solely by an AI assertion. Here ``independent''
validation or replay means a separate implementation, algorithm or library;
it does not mean independent human authorship. The implementations,
contracts, tests and their reviews were developed within the same AI
workflow. Formal proof covers the encoded properties; correctness of the
specifications, excluded implementation refinement and statistical-model
adequacy remain distinct obligations. Ghost models do not establish
correctness of foreign numerical libraries or physical I/O.

No new participant data were collected and no empirical psychometric
validation study was conducted for this project. Simulations use synthetic
draws under the stated mathematical model. The repository preserves source
and evidence history, but not a complete prompt/tool transcript, exact
per-line authorship, or an AI-contribution percentage. Those details cannot
be reconstructed reliably and have not been invented. The disclosure is
current through 1 November 2026.

\clearpage
\bibliographystyle{plainnat}
\bibliography{references}
\end{document}
